% Clase 24 --- El experimento de Fizeau (§6).
% Rueda dentada girando + espejo lejano. Lo que hay que ver: la luz sale por una
% rendija y, al volver, se encuentra con la rueda GIRADA; si el giro es justo el
% que lleva de una rendija a la siguiente, la luz pasa; si no, la tapa un diente.
\begin{center}
\begin{tikzpicture}[scale=0.95, transform shape]

  % ---- rueda dentada ----
  \foreach \ang in {0,30,...,330}{
    \fill[figgray!35] (\ang:0.55) -- (\ang:1.15)
      arc (\ang:\ang+15:1.15) -- ({\ang+15}:0.55) -- cycle;}
  \draw[figgray, line width=0.8pt] (0,0) circle (0.55);
  \draw[figgray, line width=0.8pt] (0,0) circle (1.15);
  \draw[figvec, figgray] (1.45,0.55) arc (25:65:1.55);
  \node[figlblaux, anchor=west] at (1.5,1.1) {$\omega$};
  \node[figlbl, anchor=south, align=center] at (0,1.45)
    {rueda dentada\\ (rendijas equiespaciadas)};

  % ---- fuente ----
  \draw[figvec] (-2.6,0) -- (-0.35,0);
  \node[figlbl, anchor=east] at (-2.7,0) {fuente};

  % ---- ida y vuelta ----
  \draw[figvecres] (0.35,0.14) -- (7.3,0.14);
  \draw[figvecres] (7.3,-0.14) -- (0.35,-0.14);
  \node[figlbl, figred, anchor=south] at (3.9,0.2) {ida};
  \node[figlbl, figred, anchor=north] at (3.9,-0.2) {vuelta};

  % ---- espejo ----
  \draw[figline, line width=1.8pt] (7.5,-0.9) -- (7.5,0.9);
  \foreach \yy in {-0.75,-0.45,...,0.8}{
    \draw[figgray, line width=0.6pt] (7.52,\yy) -- (7.78,{\yy+0.22});}
  \node[figlbl, anchor=north, align=center] at (7.5,-1.05) {espejo};

  % ---- distancia ----
  \draw[figgray, line width=0.6pt,
        {Latex[length=1.8mm,width=1.3mm]}-{Latex[length=1.8mm,width=1.3mm]}]
    (0,-1.9) -- (7.5,-1.9);
  \node[figlblaux, anchor=south] at (3.75,-1.86) {$L$ (dos elevaciones lejanas)};

  \node[figlbl, anchor=north, align=center] at (3.75,-2.15)
    {la luz pasa si $\dfrac{2L}{c}$ coincide con el tiempo\\
     que tarda la rueda en avanzar una rendija};

\end{tikzpicture}\\[0.5em]
{\small\itshape La dificultad de medir $c$ es que el ida y vuelta dura
poquísimo, y la salida de Fizeau fue convertir un tiempo imposible de cronometrar
en algo que sí se mide bien: una velocidad de giro. La rueda hace de obturador
periódico. A la velocidad justa, la luz que salió por una rendija encuentra al
volver la rendija siguiente ya en posición y el observador la ve; si la rueda
gira apenas más rápido o más lento, la luz regresa contra un diente y la imagen
\textbf{se apaga}. Alcanza entonces con ajustar $\omega$ hasta el apagón (o hasta
el reencendido) y conocer $L$ para despejar $c$; cuanto más grande es $L$, más
precisa la medida, y de ahí que el montaje use dos elevaciones alejadas. Que este
número, medido con espejos y engranajes, coincidiera con
$1/\sqrt{\mu_0\varepsilon_0}$ ---calculado con dos constantes de laboratorio que
nada tienen que ver con la luz--- es la razón por la que se aceptó que la luz
\textbf{es} una onda electromagnética. La $\omega$ de la rueda, dicho sea de
paso, no tiene ninguna relación con la frecuencia angular de la onda: es pura
coincidencia de símbolo.}
\end{center}
